\documentclass[12pt,a4paper]{article}

\usepackage[T1]{fontenc}
\usepackage[utf8]{inputenc}
\usepackage[brazil]{babel}
\usepackage{mathptmx}
\usepackage[a4paper,left=3cm,right=2cm,top=2.5cm,bottom=2cm]{geometry}
\usepackage{setspace}
\usepackage{indentfirst}
\usepackage{microtype}
\usepackage{hyperref}
\usepackage{xurl}

\pagestyle{empty}
\setstretch{1.12}
\setlength{\parindent}{1.25cm}
\setlength{\parskip}{0.32em}
\emergencystretch=3em
\Urlmuskip=0mu plus 3mu
\hypersetup{colorlinks=false,hidelinks}
\frenchspacing

\makeatletter
\g@addto@macro{\UrlBreaks}{\do\/\do\-\do\_\do\.\do\:}
\makeatother

\begin{document}
\sloppy

\begin{center}
{\bfseries Atividade Prática Assíncrona}\\
Geomerez Raduan de Oliveira Bandeira
\end{center}

\vspace{0.12em}
\hrule
\vspace{0.32em}

\noindent Repositório: \url{https://github.com/raduanoliveira/desenvolvimento_com_ia_aula_assincrona}.
O detalhe das etapas está em \path{ia-dev-lab/docs/relatorio-harness-etapas.md}.
O histórico desta prática está no Pull Request \url{https://github.com/raduanoliveira/desenvolvimento_com_ia_aula_assincrona/pull/3}.

\textbf{1. Autonomia e guardrail.} Executei a mesma tarefa (prazo opcional e toaster para o dia seguinte) nos modos Auto e Plan, com Composer. O Auto (\texttt{feature/prazo-auto}) concluiu em cerca de trinta minutos, com retrabalho nos testes e risco médio. O Plan (\texttt{feature/prazo-plan}) exigiu aceite do plano escrito e levou cerca de quarenta minutos; o risco percebido foi baixo. Prefiro o Auto só no tempo. Prefiro o Plan no controle e segui com essa branch. O hook em \texttt{.cursor/hooks.json} bloqueia \texttt{migrate:fresh}, \texttt{db:wipe} e a remoção do SQLite do laboratório. Na tentativa com o agente, o Cursor recusou a ação.

\textbf{2. TDD e enforcement.} No filtro \texttt{all|pending|done}, o ciclo foi teste, falha e código mínimo: Red \texttt{49ca1d1}, Green \texttt{48a55fe}, regra no Service. Sem TDD, a edição de título respeitou as camadas, mas a suíte antiga ficou verde sem cobrir o rename. Na prioridade usei o Superpowers de ponta a ponta: desenho aprovado, depois o teste que falhou e o código até voltar a verde (\texttt{0b75d75}). O plugin reforçou o \texttt{AGENTS.md}; não o substituiu.

\textbf{3. Checkpoint humano.} O transcript está em \texttt{ia-dev-lab/docs/sessao-log.md}. O ponto de parada é antes de criar ou aplicar migration. Isso não se confunde com o hook: o hook recusa destruir o banco; o checkpoint decide se o schema muda. Na simulação do Plan do \texttt{due\_date}, a decisão foi aprovar. Assumi o papel de responsável pelo schema e pelos dados do SQLite no Docker, não só de revisor de código.

\textbf{4. Decisão, ADR e diagrama.} O ToDo está no tamanho certo: extrair API de tarefas ou de login viraria acoplamento de rede sem segundo consumidor. Fechei o contrato da aplicação com \texttt{app/Domain}. O ADR 0005 registra isso no formato MADR; os ADRs 0001 a 0004 permaneceram intactos. No C4 de contêiner, a versão B comunica melhor, porque desenha o que se implanta (tela, API, SQLite, Google e GitHub), não a camada POSA.

\textbf{5. Custo e tokens.} Estimei pelos transcripts (caracteres do log divididos por 4). A prioridade consumiu cerca de 22 mil tokens; o filtro, o mesmo patamar; o prazo, cerca de 58 mil, porque foram Auto, Plan e hook. Reescrever o relatório custou mais que o filtro. A estratégia aplicável é \emph{model routing}: Composer na feature e modelo menor nos ajustes de LaTeX.

\textbf{6. Uma dificuldade real.} A dificuldade foi conceitual, no C4 de contêiner. O prompt curto promoveu Service e repositório a processo, inventou MySQL e omitiu o GitHub. Camada POSA não é unidade de implantação neste laboratório: Service e repositório vivem dentro da API. Quem lê só a versão A pode achar que existe um microsserviço de regra e outro de dados. Só o segundo prompt, dizendo o que \emph{não} é contêiner, descreveu a arquitetura do ToDo.

\end{document}
